\documentclass[10pt,a4paper]{amsart}
\usepackage[margin=19mm]{geometry}
\usepackage{amsmath,amssymb,mathtools}
\usepackage[scr=rsfs,cal=euler]{mathalfa}
\usepackage{xfrac}
\usepackage[unicode,hidelinks,bookmarks=false]{hyperref}
\newcommand{\ie}{i.e.\ }
\newcommand{\cf}{cf.\ }
\newcommand{\resp}{respectively\ }
\newcommand{\Subsection}{\subsection}
\providecommand{\nobreakdash}{\nobreak\hbox{-}}
\setlength{\emergencystretch}{2em}
\multlinegap0pt
\usepackage{bm}
\usepackage{bbm}
\usepackage{dsfont}
\usepackage{xspace}
\usepackage{cancel}
\usepackage{mathtools}
\usepackage{amsthm}
\makeatletter
\@ifundefined{theorem}{\newtheorem{theorem}{Theorem}}{}
\@ifundefined{definition}{\newtheorem{definition}[theorem]{Definition}}{}
\@ifundefined{lemma}{\newtheorem{lemma}[theorem]{Lemma}}{}
\@ifundefined{prop}{\newtheorem{prop}[theorem]{Proposition}}{}
\makeatother
\def\det{\,{\rm det\ }}
\newcommand{\R}{\mathbb{R}}
\newcommand{\res}      {\mathop{\hbox{\vrule height 7pt width .5pt depth
                        0pt\vrule height .5pt width 6pt depth 0pt}}\nolimits}
\title[On the relaxation of polyconvex functionals with linear grow]{On the relaxation of polyconvex functionals with linear growth under strict convergence in $BV$}
\author{Riccardo Scala}
\date{Source: arXiv:2508.11041v1 (CC BY 4.0)}
\begin{document}
\maketitle

\noindent\emph{Source and licence.} On the relaxation of polyconvex functionals with linear growth under strict convergence in $BV$. Version arXiv:2508.11041v1 (CC BY 4.0), distributed under the Creative Commons Attribution 4.0 International licence, as recorded in the arXiv licence metadata for this exact version and shown on the abstract page https://arxiv.org/abs/2508.11041v1. Full rights notice: licenses/arxiv-2508.11041-CC-BY-4.0.txt.\par\smallskip

\noindent\textit{Editorial scope.} Counted unit: the Lemma whose source label is \texttt{lem\_glueing1} in the article (source0.tex lines 1264--1275, source label \texttt{lem\_glueing1}), delivered together with the article's own complete original proof of it (source0.tex lines 1276--1410, one \texttt{proof} environment of 135 lines). No mathematics is written, repaired or re-derived: statement and proof are the article's own text line for line. Required context retained uncounted: growthbis (lines 524--526); gencurve (lines 611--619); prop\_convstrict (lines 627--636); Phi\_interp (lines 727--729); Psi\_interp (lines 732--734); mathcalTdelta (lines 821--823); detT-1 (lines 838--839); Sigma- (lines 1045--1047); Sigma+ (lines 1088--1090); Hinterp (lines 1108--1113); est:nablaH (lines 1135--1138); est:detH (lines 1144--1148); grad\_inTdelta (lines 1205--1209); det\_inTdelta (lines 1211--1214). Every definition, display and label of those lines is retained unchanged. Omitted from this package and not redistributed: 1--523, 527--610, 620--626, 637--726, 730--731, 735--820, 824--837, 840--1044, 1048--1087, 1091--1107, 1114--1134, 1139--1143, 1149--1204, 1210--1210, 1215--1263, 1411--2404. Nothing inside a retained excerpt is omitted or altered: every retained body is delivered exactly as the article's lines are written, so no omission receipt of any kind accompanies this package. Citations: the retained excerpts carry no citation command at all, so no bibliography entry of the article is transcribed and no reference is redistributed. Reference adaptation: none was needed and none was made; every reference inside the delivered unit resolves inside it, so no unresolved reference or citation is produced. Printed slips kept literal: no printed word, symbol, hypothesis or number of the counted statement or of its proof was altered. Layout and class: the article's own class is \texttt{article} with packages including epsf, epsfig, graphics, color, amsmath, amssymb, cite, verbatim, float, graphicx, amsthm, textcomp, subfig, hyperref; the wrapper is the lane's pinned \texttt{amsart} editorial wrapper at 10pt on A4, which restates the theorem-like environments the retained text needs and adds the article's own macro definitions that the retained text needs (\texttt{\textbackslash R}, \texttt{\textbackslash det}, \texttt{\textbackslash res}), with extra wrapper packages (bm, bbm, dsfont, xspace, cancel, mathtools). The delivered numbering is the unit's own. Assets: no retained span contains a figure environment, a graphics-inclusion command, a PDF-page inclusion, a table or a TikZ picture; no image file is copied into this package, the package declares no asset dependency and no image file is redistributed with it. Licence: version arXiv:2508.11041v1 of this article is distributed under the Creative Commons Attribution 4.0 International licence, as recorded in the arXiv licence metadata for this exact version and shown on the abstract page https://arxiv.org/abs/2508.11041v1. A full rights notice is delivered at licenses/arxiv-2508.11041-CC-BY-4.0.txt. Characters: every retained excerpt is pure ASCII, so the delivered engine, XeLaTeX with its default Latin Modern text font, reports no missing character.
\begin{align}\label{growthbis}
\|\partial g\|_{L^\infty}\leq C_g.
\end{align}

\begin{definition}
Given $\gamma\in BV([a,b];\R^m)$ we define $\overline \gamma:[0,1]\rightarrow \R^m$ as  
\begin{align}\label{gencurve}
	\overline \gamma(s)=\begin{cases}
		\frac{\gamma(t^+)(s-s_\gamma(t^-))+\gamma(t^-)(s_\gamma(t^+)-s)}{s_\gamma(t^+)-s_\gamma(t^-)}&\text{if }s\in [s_\gamma(t^-),s_\gamma(t^+)],\\
		\gamma(t_\gamma(s))&\text{otherwise.}
	\end{cases}
\end{align}
\end{definition}

\begin{prop}\label{prop_convstrict}
Let $\gamma\in BV([a,b];\R^m)$ and let $(\varphi_k)\subset \textrm{Lip}([a,b];\R^m)$ be a sequence of maps converging strictly to $\gamma$ as $k\rightarrow \infty$. The functions $\overline \varphi_k:=\varphi_k\circ t_{\varphi_k}:[0,1]\rightarrow \R^m$ are Lipschitz continuous with uniformly bounded Lipschitz constants and \begin{align}\label{conv_curves}
&\overline\varphi_k\rightarrow \overline \gamma\qquad \text{strictly in $BV([0,1];\R^m)$ and weakly star in }W^{1,\infty}([0,1];\R^m),\nonumber\\
&s_{\varphi_k}\rightarrow s_\gamma\qquad \text{strictly in }BV([a,b]),\\
&t_{\varphi_k}\rightarrow t_\gamma\qquad \text{weakly star in }W^{1,\infty}([0,1]).\nonumber
\end{align}  
Moreover there exists a function $a_\gamma:\R^+\rightarrow \R^+$ depending only on $\gamma$ and such that $a_\gamma(t)\rightarrow 0$ when $t\rightarrow0^+$, and 
$$\|s_{\varphi}-s_\gamma\|_{L^1}+\|\overline\varphi-\overline \gamma\|_{L^\infty}\leq a_\gamma(d_s(\varphi,\gamma)),$$
for all $\varphi\in \textrm{Lip}([a,b];\R^m)$. 
\end{prop}

\begin{align}\label{Phi_interp}
\Phi_{\varphi,\psi}(t,r):=\varphi\Big(t_\varphi\big(s_\varphi(t)\frac{r}{h}+s_\psi(t)\frac{h-r}{h} \big)\Big),
\end{align}

\begin{align}\label{Psi_interp}
	\Psi_{\varphi,\psi}(t,r):=\varphi\Big(t_\varphi\big(s_\psi(t)\big)\Big)\frac{h-r}{h}+\psi\Big(t_\psi\big(s_\psi(t)\big)\Big)\frac{r}{h}=\overline \varphi(s_\psi(t))\frac{h-r}{h}+\overline \psi(s_\psi(t))\frac{r}{h},
\end{align}

\begin{align}\label{mathcalTdelta}
	\mathcal T_\delta(t,r):=\gamma(t)+r\dot\gamma(t)^\perp,
\end{align}

\begin{align}\det(\nabla \mathcal T_h^{-1}(x))= 1+R'_h(x),\qquad \qquad \|R'_h\|_{L^\infty}\leq C_\gamma h.\label{detT-1}
\end{align}

\begin{align}\label{Sigma-}
\Sigma_{\delta,n,c}^-:\overline T_\delta^-\setminus T_{\frac cn}\rightarrow \overline T^-_\delta,\qquad \qquad \Sigma^-_{\delta,n,c}:=(\mathcal T_\delta \circ \Upsilon_{\delta,n,c}\circ \mathcal T_\delta^{-1})\res (\overline T^-_\delta\setminus T_{\frac cn}),
\end{align}

\begin{align}\label{Sigma+}
\Sigma^+_{\delta,n,c}:=(\widehat {\mathcal T}_\delta \circ \Upsilon_{\delta,n,c}\circ \widehat{\mathcal T}_\delta^{-1})\res (\overline T_\delta^+\setminus  T_{\frac cn}).
\end{align} 

\begin{align}\label{Hinterp}
H_{\varphi,\psi,h}:=\begin{cases}
				 \Phi_{\varphi,\psi}\circ \mathcal T_h^{-1}&\text{ in }\overline T_h^+\\
			 \Psi_{\varphi,\psi}\circ \widehat {\mathcal T}_h^{-1}&\text{ in }\overline T_h^-,
					\end{cases}
\end{align}

\begin{align}\label{est:nablaH}
	\int_{T_h}|\nabla H_{\varphi,\psi,h}(x)|dx&\leq 2C_\gamma h(L_\varphi+b-a)+C_\gamma (L_\varphi+b-a)\int_a^b|s_\psi(t)-s_\varphi(t)|dt\nonumber\\&\;\;\;+C_\gamma(L_\varphi+L_\psi+b-a)h+C_\gamma\int_0^{L_\psi}|{\overline \psi}(s)-{\overline \varphi}(s)|ds\nonumber\\
	&\leq C_{\gamma,L_\varphi,L_\psi}(h+\|s_\psi-s_\varphi\|_{L^1}+\|{\overline \psi}-{\overline \varphi}\|_{L^\infty}),
\end{align}

\begin{align}\label{est:detH}
&\int_{T_h^+}|\det(\nabla H_{\varphi,\psi,h})(x)|dx=0,\\
&\int_{T_h^-}|\det(\nabla H_{\varphi,\psi,h})(x)|dx\leq \int_D|\det\big(\nabla \Psi_{\varphi,\psi}(t,r)\big)||\det(\nabla \mathcal T_h^{-1}(\mathcal T_h(t,r)))||\det(\nabla \mathcal T_h(t,r))|dtdr\nonumber\\
&\qquad \qquad \qquad \qquad \qquad \qquad\leq (L_\varphi+L_\psi+b-a)\int_0^{1}|{\overline \psi}(s)-{\overline \varphi}(s)|ds\leq C_{\gamma,L_\varphi,L_\psi}\|{\overline \psi}-{\overline \varphi}\|_{L^\infty}.\nonumber
\end{align}

\begin{align}\label{grad_inTdelta}
	&\int_{T^-_\delta\setminus T_{\frac cn}}|\nabla u^--\nabla v|dx\leq  \int_{T^-_\delta\setminus T_{\frac cn}}|\nabla v(\Sigma_{\delta,n,c}(x))-\nabla v(x)|dx +\frac {C_\gamma}{n}\int_{T^-_\delta\setminus T_{\frac cn}}|\nabla v(\Sigma_{\delta,n,c}(x))|dx\nonumber\\
	&\qquad\qquad\qquad\qquad\qquad \leq  \beta^-_v(\frac1n)+\frac {C_\gamma}{n}(1+ \frac {C_\gamma}{n})\int_{T_\delta^-} |\nabla v|dx\nonumber\\
	&\int_{T^+_\delta\setminus T_{\frac cn}}|\nabla u^+-\nabla v|dx\leq \beta^+_v(\frac1n)+\frac {C_\gamma}{n}(1+ \frac {C_\gamma}{n})\int_{T_\delta^+} |\nabla v|dx,
\end{align}

\begin{align}\label{det_inTdelta}
	\int_{T^-_\delta\setminus T_{\frac cn}}|\det(\nabla u^-)-\det(\nabla v)|dx\leq \eta^-_v(\frac1n)+\frac {C_\gamma}{n}(1+ \frac {C_\gamma}{n})\int_{T_\delta^-}|\det(\nabla v)|dx,\nonumber\\
	\int_{T^+_\delta\setminus T_{\frac cn}}|\det(\nabla u^+)-\det(\nabla v)|dx\leq \eta^+_v(\frac1n)+\frac {C_\gamma}{n}(1+ \frac {C_\gamma}{n})\int_{T_\delta^+}|\det(\nabla v)|dx.
\end{align}

\begin{lemma}\label{lem_glueing1}
	Let $A\subset\R^2$ be a bounded  open set and let $B\subset A$ be a open subset whose boundary is $\partial B=:\Gamma\subset A$ is a  closed Jordan curve of class $C^3$; let $ u\in BV(A;\R^m)$ be such that $$|Du|(\partial B)=0,\qquad u\res\partial B\in BV(\partial B;\R^m),$$ let $v^+,v^-\in\text{{\rm Lip}}_{\text{{\rm loc}}}(A;\R^m)$ be two maps and let $\delta>0$ small so that $T_\delta$ is a tubular neighborhood of $\Gamma$. Then there exists a function $\omega_\Gamma:\R^+\rightarrow \R^+$ depending on $\Gamma$ and on $u\res \Gamma$ (but independent of $v^\pm$) with $\lim_{t\rightarrow 0^+}\omega_\Gamma(t)=0$ and the such that following holds: for all $\varepsilon>0$ there exists a function $w\in\text{{\rm Lip}}_{\text{{\rm loc}}}(A;\R^m)$ with 
	\begin{align}
	&w=v^-\text{ in }B\setminus T_\delta\text{ and }w=v^+\text{ in }A\setminus B\setminus  T_\delta,\nonumber\\
	&\|w-v^-\|_{L^1(T_\delta\cap B)}\leq 3\|v^--u\|_{L^1(T_\delta\cap B)}+r,\nonumber\\
		&\|w-v^+\|_{L^1(T_\delta\cap (A\setminus B))}\leq 3\|v^+-u\|_{L^1(T_\delta\cap (A\setminus B))}+r,\nonumber\\
		&\int_B|\nabla w-\nabla v^-|dx\leq r,\label{glueingthesis}\qquad \qquad \int_{A\setminus B}|\nabla w-\nabla v^+|dx\leq  r,\\
		&F(w, B)\leq F(v^-, B)+r,\;\; F(w,A\setminus \overline B)\leq F(v^+,A\setminus \overline B)+r,\nonumber\\
		&r\leq\varepsilon+\omega_\Gamma\big(d_s(v^+\res \Gamma,u\res\Gamma)+d_s(v^-\res \Gamma,u\res\Gamma)\big).\nonumber
	\end{align}
Moreover, if $v^+,v^-\in\text{{\rm Lip}}(A;\R^m)$ then $w\in\text{{\rm Lip}}(A;\R^m)$.
\end{lemma}

\begin{proof}
	Assume  that  $\gamma:[a,b]\rightarrow \R^2$ is an arc-length parametrization of the loop $\Gamma$. Let us consider the corresponding map $\mathcal T_\delta$ in \eqref{mathcalTdelta} and let $T^-_\delta$ and $T^+_\delta$ denote the interior and external parts of $T_\delta$ with respect to $B$, i.e., $ T^-_\delta=\overline B\cap T_\delta$, $T^+_\delta=T_\delta\setminus B$. Then we set, for any $n\geq 1$,
	\begin{align}
		w_{n}:=\begin{cases}
			v^-\circ \Sigma_{\delta,n,c}^-&\text{ in }T^-_\delta\setminus T_{\frac cn},\\
			v^+\circ  \Sigma_{\delta,n,c}^+&\text{ in }T^+_\delta\setminus T_{\frac cn},\\
			v^-&\text{ in }B\setminus T_\delta,\\
			v^+&\text{ in }A\setminus B\setminus T_\delta,
		\end{cases}
	\end{align}
	where we recall the maps $\Sigma_{\delta,n,c}^-$ and $\Sigma_{\delta,n,c}^+$ in \eqref{Sigma-} and \eqref{Sigma+}, with $c\in (0,\delta)$ fixed.
	We have to define $w_{n}$ in $T_{\frac cn}$: we set 
	$$\widetilde\varphi:=(v^-\circ \Sigma_{\delta,n,c}^- )\res \Gamma_{-\frac cn},\qquad \widetilde\psi:=(v^+\circ \Sigma_{\delta,n,c}^+)\res \Gamma_{\frac cn},$$
	and 
	recalling \eqref{Phi_interp}, \eqref{Psi_interp}, and \eqref{Hinterp}, we define
	\begin{align}
		w_{n}:=H_{\varphi,\psi,\frac cn}\qquad \qquad \text{ in }T_{\frac cn},
	\end{align}
	where $\varphi,\psi:[a,b]\rightarrow \R^m$ are given by
	\begin{align}
		\varphi=\widetilde \varphi\circ \mathcal T_\delta (\cdot,-\frac cn) \qquad \qquad \psi=\widetilde \psi\circ \mathcal T_\delta (\cdot,\frac cn).
	\end{align}
	By definition of $\widetilde\varphi$ and $\widetilde\psi$, using \eqref{Sigma-} and \eqref{Sigma+}, $\varphi$ and $\psi$  can be equivalently written as 
	\begin{align*}
		&\varphi(t)=v^-\circ\mathcal T_\delta\circ \Upsilon_{\delta,n,c}(t,-\frac cn)=v^-(\mathcal T_\delta(t,0))=v^-(\gamma(t))\\
		&\psi(t)=v^+\circ\mathcal T_\delta\circ \Upsilon_{\delta,n,c}(t,\frac cn)=v^+(\mathcal T_\delta(t,0))=v^+(\gamma(t)).
	\end{align*}
%{\color{red}
%	In particular, the hypothesis that $u_k\res\Gamma$ and $v_k\res\Gamma$ converge strictly to $u\res\Gamma$ in $BV(\Gamma;\R^m)$ implies that 
%	\begin{align}\label{conv_stricts}
%		\varphi_k\rightarrow u\circ \gamma,\qquad \psi_k\rightarrow u\circ \gamma\qquad\text{ strictly in }BV([a,b];\R^m);
%	\end{align}
%	This is due to Remark \ref{rem_stricttraces}.}
%	
	In this way we have that $w_{n}$ is Lipschitz continuous in $T_{\frac cn}$ and 
	$$w_{n}=\widetilde \varphi \text{ on }\Gamma_{-\frac cn},\qquad \qquad w_{n}=\widetilde \psi \text{ on }\Gamma_{\frac cn}.$$
	Moreover $w_n$ turns out to be globally  Lipschitz in $A$ if so are $v^+$ and $v^-$.
	Let us estimate the gradient and Jacobian determinant integral of $w_{n}$ in $T_{\delta}$: by \eqref{grad_inTdelta}  we have
	\begin{align}\label{gradient1}
	\int_{T^-_\delta\setminus T_{\frac cn}}|\nabla w_{n}-\nabla v^-|dx&\leq   \beta^-_{v^-}(\frac1n)+\frac {C_\gamma}{n}(1+ \frac {C_\gamma}{n})\int_{T_\delta^-} |\nabla v^-|dx,\nonumber\\
	\int_{T^+_\delta\setminus T_{\frac cn}}|\nabla w_{n}-\nabla v^+|dx&\leq   \beta^+_{v^+}(\frac1n)+\frac {C_\gamma}{n}(1+ \frac {C_\gamma}{n})\int_{T_\delta^+} |\nabla v^+|dx,
	\end{align}
and in particular there is a constant $C_\gamma>0$ (depending on $\gamma$, but independent of $n$) such that 
	\begin{align}
	&\int_{T_\delta\setminus T_{\frac cn}}|\nabla w_{n}|dx\leq \beta^-_{v^-}(\frac1n)+  \beta^+_{v^+}(\frac1n)+ \frac { C_\gamma}{n}\Big(\int_{T^-_\delta} |\nabla v^-|dx+\int_{T^+_\delta} |\nabla v^+|dx\Big),\label{257}
\end{align}
Furthermore, 
%	whereas \eqref{oppositeineqs} implies
%		\begin{align}\label{gradient2}
%		\int_{T^-_\delta\setminus T_{\frac cn}}|\nabla w_{n}|dx&\geq  (1- \frac {C_\gamma}{n})^2\int_{T^+_\delta} |\nabla v^-|dx,\nonumber\\
%		\int_{T^+_\delta\setminus T_{\frac cn}}|\nabla w_{n}|dx&\geq  (1- \frac {C_\gamma}{n})^2\int_{T^+_\delta} |\nabla v^+|dx.
%	\end{align}
on account of \eqref{det_inTdelta}, it follows, for all $i,j\in \{1,\dots,m\},$ $i\neq j$,
\begin{align}
&\int_{T_\delta^-\setminus T_{\frac cn}}|M_{12}^{ij}(\nabla w_{n})-M_{12}^{ij}(\nabla v^-)|dx+\int_{T_\delta^-\setminus T_{\frac cn}}|M_{12}^{ij}(\nabla w_{n})-M_{12}^{ij}(\nabla v^-)|dx\nonumber\\
&\leq  \eta^-_{v^-}(\frac1n)+ \eta^+_{v^+}(\frac1n)+{C_\gamma}{n}\Big(\int_{T^-_\delta}|M_{12}^{ij}(\nabla v^-)|dx+\int_{T^+_\delta}|M_{12}^{ij}(\nabla v^+)|dx\Big).\label{258}
\end{align}
%	\begin{align}
%		&\int_{T_\delta\setminus T_{\frac cn}}|\nabla w_{n}|dx\leq  (1+ \frac {C_\gamma}{n})^2\Big(\int_{T^-_\delta} |\nabla v^-|dx+\int_{T^+_\delta} |\nabla v^+|dx\Big),\label{257}\\
%		&\int_{T_\delta\setminus T_{\frac cn}}|\det(\nabla w_{n})|dx\leq (1+ \frac {C_\gamma}{n})\Big(\int_{T^-_\delta}|\det(\nabla v^-)|dx+\int_{T^+_\delta}|\det(\nabla v^+)|dx\Big)\label{258}
%	\end{align}
%and at the same time
%	\begin{align}
%		&\label{259}\int_{T_\delta\setminus T_{\frac cn}}|\nabla w_{n}|dx\geq  (1- \frac {C_\gamma}{n})^2\Big(\int_{T^-_\delta} |\nabla v^-|dx+\int_{T^-_\delta} |\nabla v^+|dx\Big),\\
%		&\int_{T_\delta\setminus T_{\frac cn}}|\det(\nabla w_{n})|dx\geq (1- \frac {C_\gamma}{n})\Big(\int_{T^-_\delta}|\det(\nabla v^-)|dx+\int_{T^+_\delta}|\det(\nabla v^+)|dx\Big)\label{260}
%	\end{align}
	As for the integral on  $T_{\frac cn}$, using \eqref{est:nablaH} and \eqref{est:detH},  we have for all $i,j\in \{1,\dots,m\},$ $i\neq j$,
	\begin{align}\label{261}
		\int_{T_{\frac{c}{n}}}|\nabla w_{n}|dx&\leq C_{\gamma,L_\varphi,L_\psi}(\frac{1}{n}+\|s_{\psi}-s_{\varphi}\|_{L^1}+\|{\overline \psi}-{\overline \varphi}\|_{L^\infty})\nonumber\\
		&\leq C_{\gamma,L_\varphi,L_\psi}\big(\frac1n+\|s_{\psi}-s_{\sigma}\|_{L^1}+\|s_{\varphi}-s_\sigma\|_{L^1}+\|{\overline \psi}-{\overline\sigma}\|_{L^\infty}+\|{\overline \varphi}-{\overline\sigma}\|_{L^\infty}\big)\nonumber\\
		\int_{T_{\frac{c}{n}}}|M_{12}^{ij}(\nabla w_{n})|dx&\leq C_{\gamma,L_\varphi,L_\psi}\|{\overline \psi}-{\overline \varphi}\|_{L^\infty}\leq C_{\gamma,L_\varphi,L_\psi}\big(\|{\overline \psi}-{\overline \sigma}\|_{L^\infty}+\|{\overline \varphi}-{\overline \sigma}\|_{L^\infty}\big).
	\end{align}
	Here we have set $\sigma:=u\circ \gamma$ and denoted by $\overline \sigma$ the generalized curve in \eqref{gencurve}. By Proposition \ref{prop_convstrict} we find a function $a_\gamma$ such that, up to enlarging the constant $C_{\gamma,L_\varphi,L_\psi}$ if necessary, 
	\begin{align}\label{265}
		\int_{T_{\frac{c}{n}}}|M_{12}^{ij}(\nabla w_{n})|dx+\int_{T_{\frac{c}{n}}}|\nabla w_{n}|dx\leq C_{\gamma,L_\varphi,L_\psi}\big(\frac{1}{n}+a_\gamma(d_s(\varphi,\sigma)+d_s(\psi,\sigma))\big).
	\end{align}
%	Therefore letting $n=k$, and denoting $w_k:=w_{k,k}$, to conclude the proof we first observe that the last inequality, together with \eqref{257} and \eqref{259}, implies   
%	\begin{align*}
%		\lim_{k\rightarrow \infty}\int_A |\nabla w_k|dx=\lim_{k\rightarrow \infty}\Big(\int_B |\nabla v_k|dx+\int_{A\setminus B} |\nabla u_k|dx\Big)=|Du|(B)+|Du|(A\setminus B)=|Du|(A),
%	\end{align*}
%	and since $w_k\rightarrow u$ pointwise a.e., we obtain that $w_k\rightarrow u$ strictly in $BV(A;\R^m)$.
We observe that inequalities 
	\eqref{gradient1}, \eqref{258}, and \eqref{265} entail
	\begin{align*}
		\int_{T^-_\delta}|\mathcal M(\nabla v^-)-\mathcal M(\nabla w_n)|dx+\int_{T^+_\delta}|\mathcal M(\nabla v^+)-\mathcal M(\nabla w_n)|dx
		\leq o(n)+C_\gamma a_\gamma(d_s(\varphi,\sigma)+d_s(\psi,\sigma))
	\end{align*}
	% and
	%\begin{align}
	%\int_{ T_{\frac cn}}|\mathcal M(\nabla w_k)|dx\leq \frac{C}{k},
	%\end{align}
	for some quantity $o(n)$ tending to $0$ as $n\rightarrow\infty$. These estimates together with \eqref{growthbis} entail
	\begin{align}\label{268}
		\mathcal F(w_n,B)&=\int_Bg(\mathcal M(\nabla w_n))\leq\int_Bg(\mathcal M(\nabla v^-))-\partial g(\mathcal M(\nabla v^-))(\mathcal M(\nabla v^-)-\mathcal M(\nabla w))dx\nonumber\\
		&\leq \mathcal F(v^-,B)+C_g \int_{T_\delta^-}|\mathcal M(\nabla v^-)-\mathcal M(\nabla w)| dx=:\mathcal F(v^-,B)+r',
	\end{align} 
	with $r'\leq C_g(o(n)+C_\gamma a_\gamma(d_s(\varphi,\sigma)+d_s(\psi,\sigma))) =:C_g o(n)+\omega_\Gamma(d_s(\varphi,\sigma)+d_s(\psi,\sigma))$. A similar reasoning  for the set $A\setminus B$ leads to
	\begin{align*}
	F(w,A\setminus B)=F(v^+,A\setminus B)+r'',
	\end{align*}
with $r''\leq C_g o(n)+\omega_\Gamma(d_s(\varphi,\sigma)+d_s(\psi,\sigma))$. So if we take $n$ large enough, we have obtained the last but one line in \eqref{glueingthesis}.
% Further, inequalities \eqref{257} and  \eqref{265} imply
%\begin{align*}
%\int_B|\nabla w|dx\leq \int_B|\nabla v^-|dx+\frac{\overline C}{n}+\overline Ca_\Gamma(d_s(\varphi,\sigma)+d_s(\psi,\sigma))\\
%\int_B|\nabla w|dx\leq \int_B|\nabla v^+|dx+\frac{\overline C}{n}+\overline Ca_\Gamma(d_s(\varphi,\sigma)+d_s(\psi,\sigma)),
%\end{align*}
Also  the forth inequality in \eqref{glueingthesis} easily follows from \eqref{gradient1} and \eqref{261}.
It remains to estimate the $L^1$-norms. Owing to the explicit expression of $\Phi_{\varphi,\psi}$ and $H_{\varphi,\psi,\frac cn}$ in \eqref{Phi_interp} and  \eqref{Hinterp}, denoting $h=\frac cn$, we write
\begin{align*}
\int_{T^-_{\frac{c}{n}}}|w_n|dx&\leq (1+\frac{C_\gamma}{n})\int_a^b\int_0^h|\varphi\Big(t_\varphi\big(s_\varphi(t)\frac{r}{h}+s_\psi(t)\frac{h-r}{h} \big)\Big)|drdt\\
&=(1+\frac{C_\gamma}{n})\int_a^b\int_0^h|\overline \varphi\big(s_\varphi(t)\frac{r}{h}+s_\psi(t)\frac{h-r}{h} \big)|drdt\\
&\leq \frac{c}{n}(b-a)(1+\frac{C_\gamma}{n})\big(\|\overline \varphi-\overline \sigma\|_{L^\infty}+\|\overline \sigma\|_{L^\infty}\big)\\
&\leq \frac{C_\gamma}{n} \big(a_\gamma(d_s(\varphi,\sigma))+\|\sigma\|_{L^\infty}\big),
\end{align*}
where the first inequality follows from \eqref{detT-1} and the last one from Proposition \ref{prop_convstrict}. Analogously
\begin{align*}
	\int_{T^+_{\frac{c}{n}}}|w_n|dx&\leq (1+\frac{C_\gamma}{n})\int_a^b\int_0^h|\overline \varphi(s_\psi(t))\frac{h-r}{h}+\overline \psi(s_\psi(t))\frac{r}{h}|drdt\\&
	\leq \frac{C_\gamma}{n}\big(\|\overline \varphi-\overline \sigma\|_{L^\infty}+\|\overline \psi-\overline \sigma\|_{L^\infty}+2\|\overline \sigma\|_{L^\infty}\big)\\
	&\leq \frac{C_\gamma}{n} \big(a_\gamma(d_s(\varphi,\sigma))+a_\gamma(d_s(\psi,\sigma))+\|\sigma\|_{L^\infty}\big).
\end{align*}
At the same time we have
\begin{align*}
\int_B&|w_n-v^-|dx=\int_{T_\delta\setminus T_{\frac cn}}|v^-\circ\Sigma_{\delta,n,c}^- -v^-|dx\\
&\leq  \int_{T_\delta\setminus T_{\frac cn}}|v^-\circ\Sigma_{\delta,n,c}^- -u\circ\Sigma_{\delta,n,c}^-|dx+\int_{T_\delta\setminus T_{\frac cn}}|u\circ\Sigma_{\delta,n,c}^--u|dx+\int_{T_\delta\setminus T_{\frac cn}}|u-v^-|dx\\
&\leq (1+\frac{C_\gamma}{n})\|v^--u\|_{L^1(T_\delta^-)}+\int_{T_\delta\setminus T_{\frac cn}}|u\circ\Sigma_{\delta,n,c}^--u|dx+\|v^--u\|_{L^1(T_\delta^-)}\\
&\leq 3 \|v^--u\|_{L^1(T_\delta^-)}+\int_{T_\delta}|u\circ\Sigma_{\delta,n,c}-u|dx,
\end{align*}
(for $n>1/C_\gamma$)
and a similar inequality holds for $\int_{A\setminus B}|w_n-v^+|dx$.
Hence, the second and third inequalities in \eqref{glueingthesis} follow from the last three expressions, noticing that we can choose $n$ big enough so that 
\begin{align*}
\frac {C_\gamma}{n}\|\sigma\|_{L^\infty}\leq \varepsilon,\qquad \qquad \int_{T_\delta}|u\circ\Sigma_{\delta,n,c}-u|dx\leq \varepsilon,
\end{align*}
where the last condition can be obtained because  $u\circ\Sigma_{\delta,n,c}\rightarrow u$ in $L^1(T_\delta;\R^m)$ as $n\rightarrow \infty$.
\end{proof}

\end{document}
